\documentclass[11pt,a4paper]{article}

\usepackage[a4paper,margin=2.4cm]{geometry}
\usepackage{amsmath,amssymb}
\usepackage{booktabs}
\usepackage{longtable}
\usepackage{hyperref}
\usepackage{enumitem}
\usepackage[T1]{fontenc}
\usepackage{lmodern}

\hypersetup{
  colorlinks=true,
  linkcolor=blue,
  urlcolor=blue,
  pdftitle={Manual do Projeto BH\_Torus},
  pdfauthor={BH\_Torus}
}

\setlist[itemize]{noitemsep,topsep=3pt}
\setlist[enumerate]{noitemsep,topsep=3pt}
\emergencystretch=2em

\title{Manual do Projeto BH\_Torus\\
\large Collapsar, Neutrinos UHE, Toro de Matéria e Ray Tracing em Kerr}
\author{Projeto BH\_Torus}
\date{\today}

\begin{document}

\maketitle
\tableofcontents
\newpage

\section{Visão geral física e computacional}

Este projeto calcula imagens sintéticas de neutrinos de energia ultra-alta
(UHE) produzidos em uma fonte circular interna de um collapsar e propagados
através de um toro denso em torno de um buraco negro de Kerr. A simulação
combina:

\begin{itemize}
  \item geodésicas nulas em uma métrica de Kerr;
  \item uma câmera distante em coordenadas de Boyer-Lindquist;
  \item um perfil paramétrico de densidade, temperatura e $Y_e$ para o toro;
  \item seção de choque neutrino-núcleon de corrente carregada, lida de tabela;
  \item profundidade óptica DIS ao longo dos raios;
  \item emissividade UHE local e transferência radiativa atenuada;
  \item um segundo canal térmico MeV com emissividade, opacidades,
  leakage e estimativa de neutrinosfera.
\end{itemize}

O fluxo principal do programa é:

\begin{enumerate}
  \item gerar e salvar as geodésicas da câmera em
  \texttt{output/rays/kerr\_geodesics.bin};
  \item reler esse cache para cada energia observada $E_{\nu,\rm obs}$;
  \item integrar a emissão e a absorção ao longo de cada raio;
  \item salvar a imagem em \texttt{output/images/};
  \item gerar figuras em \texttt{plots/}.
\end{enumerate}

\subsection*{Status dos backgrounds de densidade}

Os perfis de densidade usados neste projeto devem ser interpretados como
\emph{morfologias semi-analíticas controladas} para estudos de opacidade UHE.
Eles não são soluções hidrodinâmicas, simulações realistas de collapsar, nem
discos de acreção em equilíbrio. O objetivo é isolar como diferentes colunas de
matéria parametrizadas afetam a profundidade óptica DIS.

A validação CUDA dos novos perfis está pendente. Antes de usar o backend CUDA
para estes backgrounds no paper, deve-se comparar explicitamente
$\tau_{\rm CPU}$ contra $\tau_{\rm CUDA}$ e $P_{\rm surv,CPU}$ contra
$P_{\rm surv,CUDA}$ para pelo menos um caso gaussiano e um caso
\texttt{powerlaw\_funnel\_envelope}.

Por padrão, o \texttt{Makefile} usa o alvo \texttt{kerr-criar-cache}, que
executa as duas etapas: geodésicas e imagem.

\section{Unidades e constantes}

O código usa unidades geométricas para a parte relativística, com

\begin{equation}
  r_g = \frac{GM}{c^2}.
\end{equation}

As coordenadas dos raios são armazenadas em unidades de $r_g$. Quando a
profundidade óptica é integrada, o elemento de caminho adimensional
\texttt{dl\_rg} é convertido para centímetros por

\begin{equation}
  dl_{\rm cm} = dl_{r_g}\,r_g.
\end{equation}

A densidade bariônica é aproximada por

\begin{equation}
  n_b = \frac{\rho}{m_u},
\end{equation}

onde $m_u$ é a unidade de massa atômica.

\subsection*{Onde está no programa}

\begin{itemize}
  \item \texttt{include/constants.hpp}: define $c$, $G$, $M_\odot$, $m_u$ e
  a função \texttt{rg\_cm(M\_msun)}.
  \item \texttt{src/radiative\_transfer.cpp}: converte \texttt{dl\_rg} para
  \texttt{dl\_cm}.
  \item \texttt{src/optical\_depth.cpp}: usa $n_b=\rho/m_u$ na profundidade
  óptica de teste para raios retos.
\end{itemize}

\section{Métrica de Kerr}

A métrica é implementada em coordenadas de Boyer-Lindquist
$(t,r,\theta,\phi)$, com assinatura $(-,+,+,+)$ e spin adimensional $a$.
As funções auxiliares são

\begin{align}
  \Sigma &= r^2 + a^2\cos^2\theta,\\
  \Delta &= r^2 - 2r + a^2,\\
  A &= (r^2+a^2)^2 - a^2\Delta\sin^2\theta.
\end{align}

O horizonte externo é

\begin{equation}
  r_+ = 1+\sqrt{1-a^2}.
\end{equation}

Os componentes covariantes usados são

\begin{align}
  g_{tt} &= -\left(1-\frac{2r}{\Sigma}\right),\\
  g_{t\phi} &= -\frac{2ar\sin^2\theta}{\Sigma},\\
  g_{rr} &= \frac{\Sigma}{\Delta},\\
  g_{\theta\theta} &= \Sigma,\\
  g_{\phi\phi} &=
  \left(r^2+a^2+\frac{2a^2r\sin^2\theta}{\Sigma}\right)\sin^2\theta.
\end{align}

O código também usa a métrica inversa

\begin{align}
  g^{tt} &= -\frac{A}{\Sigma\Delta},\\
  g^{t\phi} &= -\frac{2ar}{\Sigma\Delta},\\
  g^{rr} &= \frac{\Delta}{\Sigma},\\
  g^{\theta\theta} &= \frac{1}{\Sigma},\\
  g^{\phi\phi} &= \frac{\Delta-a^2\sin^2\theta}
  {\Sigma\Delta\sin^2\theta}.
\end{align}

Para construir a câmera local, o programa usa ainda o lapse e a velocidade
angular de arrasto de referenciais,

\begin{align}
  \alpha &= \sqrt{\frac{\Sigma\Delta}{A}},\\
  \omega &= \frac{2ar}{A}.
\end{align}

\subsection*{Onde está no programa}

\begin{itemize}
  \item \texttt{src/kerr\_metric.cpp}: \texttt{Sigma}, \texttt{Delta},
  \texttt{A}, \texttt{metric}, \texttt{inverse\_metric},
  \texttt{horizon\_radius}, \texttt{lapse} e
  \texttt{omega\_frame\_drag}.
  \item \texttt{include/kerr\_metric.hpp}: interface da classe
  \texttt{KerrMetric}.
  \item \texttt{apps/test\_kerr\_metric.cpp}: teste básico da métrica.
\end{itemize}

\section{Geodésicas nulas por Hamiltoniano}

Cada raio é propagado como uma geodésica nula. O estado canônico é

\begin{equation}
  y = (t,r,\theta,\phi,p_t,p_r,p_\theta,p_\phi).
\end{equation}

O Hamiltoniano usado é

\begin{equation}
  H = \frac{1}{2}g^{\mu\nu}p_\mu p_\nu.
\end{equation}

Para uma geodésica nula ideal, $H=0$. As equações de Hamilton são

\begin{align}
  \frac{dx^\mu}{d\lambda} &= \frac{\partial H}{\partial p_\mu}
  = g^{\mu\nu}p_\nu,\\
  \frac{dp_\mu}{d\lambda} &= -\frac{\partial H}{\partial x^\mu}
  = -\frac{1}{2}\frac{\partial g^{\alpha\beta}}
  {\partial x^\mu}p_\alpha p_\beta.
\end{align}

Como a métrica é estacionária e axissimétrica, o código mantém

\begin{equation}
  \frac{dp_t}{d\lambda}=0,\qquad
  \frac{dp_\phi}{d\lambda}=0.
\end{equation}

As derivadas espaciais da métrica inversa em $r$ e $\theta$ são avaliadas por
diferenças finitas centradas. A integração usa um passo adaptativo tipo
Runge-Kutta-Fehlberg, com fallback para RK4.

\subsection*{Onde está no programa}

\begin{itemize}
  \item \texttt{src/kerr\_geodesic.cpp}: \texttt{hamiltonian},
  \texttt{rhs}, \texttt{dg\_inv\_dr}, \texttt{dg\_inv\_dtheta},
  \texttt{step\_adaptive} e \texttt{step\_rk4}.
  \item \texttt{include/geodesic\_state.hpp}: estrutura do estado canônico.
  \item \texttt{include/kerr\_geodesic.hpp}: interface do integrador.
\end{itemize}

\section{Câmera de Kerr e condição inicial dos raios}

Cada pixel $(i,j)$ é convertido em coordenadas angulares de câmera:

\begin{align}
  u &= \left[2\frac{i+1/2}{N_x}-1\right]\tan\frac{\mathrm{FOV}}{2},\\
  v &= \left[2\frac{j+1/2}{N_y}-1\right]\tan\frac{\mathrm{FOV}}{2}.
\end{align}

A direção local normalizada é

\begin{align}
  n_r &= -\frac{1}{\sqrt{1+u^2+v^2}},\\
  n_\theta &= \frac{v}{\sqrt{1+u^2+v^2}},\\
  n_\phi &= \frac{u}{\sqrt{1+u^2+v^2}}.
\end{align}

O sinal negativo em $n_r$ indica que o raio é lançado para trás a partir do
observador, isto é, para dentro em direção ao buraco negro e ao toro.

Na aproximação tetrádica ZAMO/LNRF usada pelo código, os componentes
contravariantes iniciais são

\begin{align}
  p^t &= \frac{1}{\alpha},\\
  p^r &= \frac{n_r}{\sqrt{g_{rr}}},\\
  p^\theta &= \frac{n_\theta}{\sqrt{g_{\theta\theta}}},\\
  p^\phi &= \frac{n_\phi}{\sqrt{g_{\phi\phi}}}+\omega p^t.
\end{align}

O programa então baixa o índice via

\begin{equation}
  p_\mu = g_{\mu\nu}p^\nu.
\end{equation}

O raio é interrompido se cruza $r_+ + 10^{-3}$ ou se escapa para
\texttt{CAM\_R\_MAX\_RG}. O cache salva os pontos do caminho e uma flag
\texttt{captured}.

\subsection*{Onde está no programa}

\begin{itemize}
  \item \texttt{src/kerr\_camera.cpp}: \texttt{initial\_state} constrói
  $u,v,n_i,p^\mu,p_\mu$; \texttt{trace\_pixel} integra o raio e salva
  pontos.
  \item \texttt{include/ray.hpp}: define \texttt{PathPoint} e
  \texttt{RayPath}.
  \item \texttt{apps/compute\_kerr\_geodesics.cpp}: gera o arquivo binário
  \texttt{output/rays/kerr\_geodesics.bin}.
\end{itemize}

\section{Modelo paramétrico do toro do collapsar}

O toro é representado por um perfil gaussiano em raio e em afastamento
angular do plano equatorial. Com

\begin{equation}
  \delta_\theta = \theta-\frac{\pi}{2},
\end{equation}

a densidade é

\begin{equation}
  \rho(r,\theta) =
  \rho_0
  \exp\left[-\left(\frac{r-r_0}{\sigma_r}\right)^2\right]
  \exp\left[-\left(\frac{\delta_\theta}{H/R}\right)^2\right],
\end{equation}

mas somente dentro da região lógica atualmente fixada por

\begin{equation}
  4 < r/r_g < 18,\qquad |\theta-\pi/2| < 0.45.
\end{equation}

A temperatura parametrizada é

\begin{equation}
  T(r,\theta) =
  6\,{\rm MeV}\,
  \left[
  \exp\left[-\left(\frac{r-r_0}{\sigma_r}\right)^2\right]
  \right]^{0.2}
  \left[
  \exp\left[-\left(\frac{\delta_\theta}{H/R}\right)^2\right]
  \right]^{0.1},
\end{equation}

com piso numérico $10^{-10}$ MeV. A fração eletrônica usada é

\begin{equation}
  Y_e(r) = 0.2 + 0.1
  \exp\left[-\left(\frac{r-r_0}{4}\right)^2\right],
\end{equation}

dentro do toro.

\subsection*{Onde está no programa}

\begin{itemize}
  \item \texttt{src/torus\_profile.cpp}: \texttt{in\_torus},
  \texttt{rho}, \texttt{temperature\_MeV} e \texttt{Ye}.
  \item \texttt{include/torus\_profile.hpp}: parâmetros padrão da classe.
  \item \texttt{apps/dump\_torus\_grid.cpp}: exporta uma grade do toro.
\end{itemize}

\section{Escala física da densidade do toro}

O parâmetro \texttt{TORUS\_RHO0} deve ser interpretado como a densidade
máxima ou central do perfil paramétrico do toro, isto é, aproximadamente
$\rho_{\rm max}$ no plano equatorial perto de $r=r_0$. Ele não é a densidade
média do disco, nem a densidade do funil polar. Essa distinção é importante,
pois em um collapsar real a densidade varia por muitas ordens de grandeza
entre o disco interno, o toro externo, as regiões de fallback e o canal polar.

A literatura sugere as seguintes escalas de ordem de grandeza:

\begin{longtable}{p{4.4cm}p{4.0cm}p{6.0cm}}
\toprule
Cenário físico & Densidade típica & Interpretação para este modelo\\
\midrule
\endhead
Canal polar evacuado em simulações de collapsar &
$\rho \sim 10^6\ {\rm g\,cm^{-3}}$ &
Representa material de baixa densidade fora do toro equatorial; não deve ser
usado como densidade central de um toro denso, mas é útil como limite baixo
para testes de transparência \cite{MacFadyenWoosley1999Web}.\\
Disco/toro global em simulações hidrodinâmicas de collapsar, a alguns
segundos após o colapso &
$\rho \sim 10^9\ {\rm g\,cm^{-3}}$ no disco denso, com densidades maiores
esperadas mais internamente &
Bom valor inicial para representar um toro equatorial global sem assumir
imediatamente um NDAF extremamente compacto \cite{MacFadyenWoosley1999,
MacFadyenWoosley1999Web}.\\
Tori externos/tardios associados ao fallback, com queima nuclear no disco &
$\rho_{\rm max}\sim 5\times 10^5$--$5\times 10^8\ {\rm g\,cm^{-3}}$ &
Relevante para fases mais tardias, raios de circularização maiores
($\sim 10^8$--$10^9$ cm) e discos menos compactos
\cite{ZenatiSiegelMetzger2020}.\\
Disco interno de hiperdecreção resfriado por neutrinos, ou NDAF &
$\rho \sim 10^{10}$--$10^{13}\ {\rm g\,cm^{-3}}$ &
Regime mais denso e compacto, apropriado perto do buraco negro quando a taxa
de acreção é alta e o resfriamento por neutrinos domina \cite{LiuGuZhang2017}.\\
NDAF compacto com acoplamento magnético buraco negro--disco &
$\rho_{\rm max}\sim 10^{13}\ {\rm g\,cm^{-3}}$ em soluções de alta acreção &
Limite alto plausível para varreduras paramétricas densas, mas com física
mais extrema que o toro gaussiano simples usado aqui \cite{SheLiuXue2022}.\\
\bottomrule
\end{longtable}

Assim, para imagens e diagnósticos publicáveis, uma varredura razoável para
o parâmetro central do toro é

\begin{equation}
  \rho_0 =
  10^6,\ 10^7,\ 10^8,\ 10^9,\ 10^{10},\ 10^{11}
  \ {\rm g\,cm^{-3}},
\end{equation}

com extensão para $10^{12}$--$10^{13}\ {\rm g\,cm^{-3}}$ caso se deseje
explorar o limite NDAF interno. Valores como $\rho_0=10^{-2}\ {\rm g\,cm^{-3}}$
devem ser entendidos apenas como testes numéricos de código e organização de
saída, não como densidades físicas de um toro de collapsar.

Como a profundidade óptica DIS usada no código escala aproximadamente como

\begin{equation}
  \tau_{\rm DIS}
  \simeq
  \int \frac{\rho}{m_u}\,
  \sigma_{\nu N}^{\rm CC}(E_\nu)\,dl,
\end{equation}

a mudança de \texttt{TORUS\_RHO0} altera diretamente a atenuação dos neutrinos
UHE. Em primeira aproximação, mantendo geometria, energia e trajetória fixas,
tem-se $\tau_{\rm DIS}\propto\rho_0$. Portanto, a varredura em densidade é
fisicamente tão importante quanto a varredura em energia.

\section{Seção de choque DIS neutrino-núcleon}

A absorção usa seção de choque de corrente carregada
$\sigma_{\nu N}^{\rm CC}(E_\nu)$ lida de tabela. O arquivo padrão é

\begin{equation}
  \texttt{data/sigma/sigma\_nuN\_CC\_GBW.dat}.
\end{equation}

A tabela contém energia em GeV, seção de choque em GeV$^{-2}$ e seção de
choque em cm$^2$. O código usa a coluna em cm$^2$ para a profundidade óptica.
A interpolação é log-log:

\begin{equation}
  \ln \sigma(E) =
  \ln\sigma_1 +
  \frac{\ln E-\ln E_1}{\ln E_2-\ln E_1}
  \left(\ln\sigma_2-\ln\sigma_1\right).
\end{equation}

\subsection*{Onde está no programa}

\begin{itemize}
  \item \texttt{src/sigma\_table.cpp}: leitura da tabela e interpolação
  log-log em \texttt{log\_interp}.
  \item \texttt{include/sigma\_table.hpp}: interface da tabela.
  \item \texttt{apps/compute\_kerr\_image\_from\_cache.cpp}: escolhe
  \texttt{sigma\_nuN\_CC\_GBW.dat}.
  \item \texttt{data/sigma/}: contém modelos \texttt{GBW} e \texttt{IIM}.
\end{itemize}

\section{Profundidade óptica e probabilidade de sobrevivência}

Ao longo de um raio, a profundidade óptica DIS é

\begin{equation}
  \tau(E_\nu) =
  \int n_b(s)\,
  \sigma_{\nu N}^{\rm CC}\!\left(E_{\nu,\rm local}(s)\right)\,ds.
\end{equation}

Na forma discretizada usada no programa,

\begin{equation}
  \Delta\tau_k =
  \frac{\rho_k}{m_u}\,
  \sigma_{\nu N}^{\rm CC}(E_{\nu,\rm local,k})\,
  \Delta l_{{\rm cm},k},
\end{equation}

e

\begin{equation}
  \tau = \sum_k \Delta\tau_k.
\end{equation}

A probabilidade de sobrevivência é

\begin{equation}
  P_{\rm surv} = e^{-\tau}.
\end{equation}

O fator de redshift armazenado no caminho é usado para converter energia
observada em energia local. No integrador principal de transferência
radiativa, o código define

\begin{equation}
  g = \frac{1}{\texttt{redshift\_factor}},\qquad
  E_{\nu,\rm local} = \frac{E_{\nu,\rm obs}}{g}.
\end{equation}

Na implementação atual da câmera, \texttt{redshift\_factor} é aproximado por

\begin{equation}
  \texttt{redshift\_factor} =
  \left(1-\frac{2}{r}\right)^{-1/2},
\end{equation}

com proteção numérica perto do horizonte.

\subsection*{Onde está no programa}

\begin{itemize}
  \item \texttt{src/radiative\_transfer.cpp}: integração principal de
  $\tau$, $P_{\rm surv}$ e $I_{\rm obs}$ em \texttt{integrate\_kerr\_ray}.
  \item \texttt{src/kerr\_camera.cpp}: armazena \texttt{redshift\_factor}
  em cada \texttt{PathPoint}.
  \item \texttt{src/optical\_depth.cpp}: contém funções auxiliares para
  raios retos e para integração ao longo de \texttt{RayPath}.
\end{itemize}

\section{Fonte circular UHE e transferência radiativa}

A fonte interna é um anel gaussiano no raio e na latitude, com espectro de
potência e corte exponencial. A emissividade local é

\begin{equation}
  j(E_{\nu,\rm local},r,\theta) =
  j_0
  \exp\left[-\left(\frac{r-r_s}{\sigma_s}\right)^2\right]
  \exp\left[-\left(\frac{\theta-\pi/2}{\theta_s}\right)^2\right]
  E_{\nu,\rm local}^{-p}
  \exp\left(-\frac{E_{\nu,\rm local}}{E_{\rm max}}\right).
\end{equation}

A intensidade observada acumulada no pixel é discretizada como

\begin{equation}
  I_{\rm obs} =
  \sum_k
  g_k^3\,
  j_k\,
  e^{-\tau_k}\,
  \Delta l_{{\rm cm},k}.
\end{equation}

O fator $g^3$ representa a transformação relativística usual da intensidade
específica, compatível com a invariância de $I_\nu/\nu^3$.

\subsection*{Onde está no programa}

\begin{itemize}
  \item \texttt{src/radiative\_transfer.cpp}: função
  \texttt{emissivity\_collapsar\_ring} e soma de $I_{\rm obs}$.
  \item \texttt{include/radiative\_transfer.hpp}: define
  \texttt{UHECircularSourceParams}.
  \item \texttt{apps/compute\_kerr\_image\_from\_cache.cpp}: lê os
  parâmetros da fonte a partir dos argumentos vindos do \texttt{Makefile}.
\end{itemize}

\section{Canal térmico MeV, opacidades e leakage}

Além do canal UHE vindo do anel central, o código integra um canal térmico
MeV associado ao toro quente. Esse canal usa a mesma geodésica de Kerr, mas
outras emissividades e opacidades. A energia observada padrão é
\texttt{MEV\_ENU} em MeV.

A forma espectral local é aproximada por uma distribuição de Fermi-Dirac sem
potencial químico explícito:

\begin{equation}
  f_E(E,T) = \frac{E^2}{\exp(E/T)+1}.
\end{equation}

A emissividade térmica implementada é uma parametrização local:

\begin{equation}
  j_{\rm MeV} =
  C_{\rm MeV}
  \left[
  n_b X_{\rm cc} T^6
  + 0.05\,n_b T^9
  \right]
  f_E(E_{\nu,\rm local},T),
\end{equation}

onde

\begin{equation}
  X_{\rm cc} = (1-Y_e)+\frac{1}{2}Y_e.
\end{equation}

O primeiro termo representa emissão tipo captura carregada em núcleons; o
segundo é um termo fenomenológico tipo pares, útil para manter uma fonte
térmica sensível a temperaturas altas. A normalização $C_{\rm MeV}$ é
\texttt{MEV\_NORM}.

A opacidade de absorção efetiva é

\begin{equation}
  \kappa_{\rm abs} =
  n_b\,\sigma_{\rm abs,0}
  \left(\frac{E_{\nu,\rm local}}{\rm MeV}\right)^2
  \left[(1-Y_e)+\frac{1}{4}Y_e\right],
\end{equation}

com $\sigma_{\rm abs,0}=9.6\times10^{-44}\,{\rm cm^2}$. A opacidade de
espalhamento elástico em núcleons é

\begin{equation}
  \kappa_{\rm sca} =
  n_b\,\sigma_{\rm sca,0}
  \left(\frac{E_{\nu,\rm local}}{\rm MeV}\right)^2
  \left(1+\frac{1}{2}Y_e\right),
\end{equation}

com $\sigma_{\rm sca,0}=1.7\times10^{-44}\,{\rm cm^2}$. A profundidade óptica
MeV é

\begin{equation}
  \tau_{\rm MeV} =
  \int (\kappa_{\rm abs}+\kappa_{\rm sca})\,dl.
\end{equation}

Para representar aprisionamento difusivo sem resolver transporte de Boltzmann,
o código usa um fator leakage local,

\begin{equation}
  \Lambda(\tau_{\rm MeV}) =
  \frac{1}{1+\tau_{\rm MeV}^2},
\end{equation}

e acumula a intensidade térmica como

\begin{equation}
  I_{\rm MeV,obs} =
  \sum_k
  g_k^3\,
  j_{{\rm MeV},k}\,
  \Lambda_k\,
  e^{-\tau_{{\rm MeV},k}}\,
  \Delta l_{{\rm cm},k}.
\end{equation}

A neutrinosfera é estimada no primeiro ponto do raio, visto a partir do
observador para dentro, onde

\begin{equation}
  \tau_{\rm MeV} \geq \frac{2}{3}.
\end{equation}

\subsection*{Onde está no programa}

\begin{itemize}
  \item \texttt{src/radiative\_transfer.cpp}: funções
  \texttt{emissivity\_mev\_thermal},
  \texttt{opacity\_mev\_absorption\_cm\_inv},
  \texttt{opacity\_mev\_scattering\_cm\_inv} e acumulação de
  $I_{\rm MeV,obs}$.
  \item \texttt{include/radiative\_transfer.hpp}: define
  \texttt{MeVThermalParams} e os campos MeV de \texttt{RTResult}.
  \item \texttt{src/torus\_profile.cpp}: fornece $\rho$, $T$ e $Y_e$.
  \item \texttt{apps/compute\_kerr\_image\_from\_cache.cpp}: escreve as
  colunas MeV no arquivo de imagem.
\end{itemize}

\section{Arquivos de saída}

\subsection{Cache de geodésicas}

O arquivo binário

\begin{equation}
  \texttt{output/rays/kerr\_geodesics.bin}
\end{equation}

contém uma assinatura \texttt{KGEO}, versão, tamanho da câmera, spin usado e
os raios de cada pixel. Para cada raio, salva:

\begin{itemize}
  \item índice do pixel;
  \item coordenadas de câmera \texttt{alpha\_rg} e \texttt{beta\_rg};
  \item flag de captura;
  \item lista de pontos \texttt{PathPoint}.
\end{itemize}

\subsection{Imagem de neutrinos}

A saída de imagem é texto ASCII:

\begin{equation}
  \texttt{output/images/kerr\_image\_Enu\_<energia>.dat}.
\end{equation}

As colunas são:

\begin{center}
\begin{tabular}{ll}
\toprule
Coluna & Significado\\
\midrule
\texttt{i}, \texttt{j} & pixel da câmera\\
\texttt{alpha}, \texttt{beta} & coordenadas angulares do pixel\\
\texttt{tau\_uhe} & profundidade óptica DIS UHE integrada\\
\texttt{P\_surv\_uhe} & probabilidade UHE de sobrevivência $e^{-\tau}$\\
\texttt{I\_obs\_uhe} & intensidade UHE observada acumulada\\
\texttt{captured} & indica se o raio caiu no horizonte\\
\texttt{tau\_mev} & profundidade óptica térmica MeV\\
\texttt{P\_surv\_mev} & probabilidade MeV de sobrevivência\\
\texttt{I\_obs\_mev} & intensidade térmica MeV observada\\
\texttt{r\_neutrinosphere\_rg} & raio onde $\tau_{\rm MeV}\geq2/3$, ou -1\\
\texttt{leakage\_mev} & leakage médio ponderado pela emissividade\\
\bottomrule
\end{tabular}
\end{center}

\section{Uso do programa pelo Makefile}

Os parâmetros principais ficam no começo do \texttt{Makefile} e podem ser
alterados na linha de comando. O alvo padrão é:

\begin{verbatim}
make
\end{verbatim}

equivalente a:

\begin{verbatim}
make kerr-criar-cache
\end{verbatim}

\subsection{Rodar a imagem Kerr completa}

\begin{verbatim}
make kerr-criar-cache
\end{verbatim}

Esse comando:

\begin{enumerate}
  \item cria diretórios de saída;
  \item compila os executáveis necessários;
  \item gera o cache de geodésicas;
  \item calcula a imagem UHE para \texttt{ENU} e o canal térmico para
  \texttt{MEV\_ENU};
  \item roda \texttt{scripts/plot\_kerr\_image.py}.
\end{enumerate}

Exemplo com energia, spin e resolução:

\begin{verbatim}
make kerr-criar-cache ENU=1e10 ASPIN=0.9 CAM_NX=200 CAM_NY=200
\end{verbatim}

\subsection{Gerar somente o cache de geodésicas}

\begin{verbatim}
make kerr-geodesics ASPIN=0.9 CAM_NX=100 CAM_NY=100
\end{verbatim}

Use esse alvo quando a geometria da câmera mudou. O cache depende de
\texttt{ASPIN}, \texttt{CAM\_R\_OBS\_RG}, \texttt{CAM\_THETA\_DEG},
\texttt{CAM\_FOV\_DEG}, \texttt{CAM\_NX}, \texttt{CAM\_NY},
\texttt{CAM\_R\_MAX\_RG} e \texttt{CAM\_STEP}.

\subsection{Reusar o cache para outra energia}

\begin{verbatim}
make image-from-cache ENU=1e12
\end{verbatim}

Use esse alvo quando a câmera e o spin não mudaram, mas a energia ou os
parâmetros do toro/fonte mudaram. Isso evita recalcular as geodésicas.

\subsection{Plotar uma imagem já calculada}

\begin{verbatim}
make plot-kerr-image ENU=1e9
\end{verbatim}

Esse alvo chama:

\begin{verbatim}
python3 scripts/plot_kerr_image.py 1e9
\end{verbatim}

\subsection{Visualizar o toro}

\begin{verbatim}
make plot-torus
\end{verbatim}

Esse alvo gera a grade do toro e chama o script de visualização.

\subsection{Teste de profundidade óptica em raio reto}

\begin{verbatim}
make tau-straight ENU=1e9
\end{verbatim}

Esse fluxo usa a aproximação auxiliar de raios retos em
\texttt{src/optical\_depth.cpp}.

\subsection{Testes}

\begin{verbatim}
make tests
\end{verbatim}

Atualmente o alvo compila e roda o teste da métrica de Kerr.

\subsection{Limpeza}

\begin{verbatim}
make clean-build
make clean-output
make clean
\end{verbatim}

\section{Parâmetros principais do Makefile}

\begin{longtable}{llp{7.3cm}}
\toprule
Parâmetro & Padrão & Significado\\
\midrule
\endhead
\texttt{ENU} & \texttt{1e9} & Energia observada do neutrino em GeV.\\
\texttt{ASPIN} & \texttt{0.0001} & Spin adimensional do buraco negro.\\
\texttt{MBH\_MSUN} & \texttt{3.0} & Massa do buraco negro em massas solares.\\
\texttt{CAM\_R\_OBS\_RG} & \texttt{60.0} & Raio do observador em $r_g$.\\
\texttt{CAM\_THETA\_DEG} & \texttt{72.0} & Inclinação polar do observador em graus.\\
\texttt{CAM\_FOV\_DEG} & \texttt{20.0} & Campo de visão da câmera.\\
\texttt{CAM\_NX}, \texttt{CAM\_NY} & \texttt{100}, \texttt{100} & Resolução da câmera.\\
\texttt{CAM\_R\_MAX\_RG} & \texttt{120.0} & Raio externo para parar raios que escapam.\\
\texttt{CAM\_STEP} & \texttt{0.001} & Passo base do integrador de geodésicas.\\
\texttt{TORUS\_RHO0} & \texttt{1.0e-2} & Densidade central do toro em g cm$^{-3}$.\\
\texttt{TORUS\_R0\_RG} & \texttt{10.0} & Raio central do toro em $r_g$.\\
\texttt{TORUS\_SIGMA\_RG} & \texttt{5.0} & Largura radial gaussiana do toro.\\
\texttt{TORUS\_H\_OVER\_R} & \texttt{0.25} & Espessura angular vertical aproximada.\\
\texttt{SOURCE\_R\_RG} & \texttt{3.5} & Raio central da fonte circular UHE.\\
\texttt{SOURCE\_SIGMA\_RG} & \texttt{1.0} & Largura radial da fonte.\\
\texttt{SOURCE\_THETA\_DEG} & \texttt{15.0} & Espessura angular da fonte em graus.\\
\texttt{SOURCE\_POWERLAW} & \texttt{2.0} & Índice espectral $p$ em $E^{-p}$.\\
\texttt{SOURCE\_EMAX\_GEV} & \texttt{1.0e12} & Energia de corte exponencial.\\
\texttt{SOURCE\_NORM} & \texttt{1.0} & Normalização arbitrária da emissividade.\\
\texttt{MEV\_ENU} & \texttt{10.0} & Energia observada do canal térmico MeV.\\
\texttt{MEV\_NORM} & \texttt{1.0} & Normalização arbitrária da emissividade MeV.\\
\texttt{OMP\_NUM\_THREADS} & \texttt{6} & Número de threads OpenMP.\\
\bottomrule
\end{longtable}

\section{Mapa rápido: teoria, código e comandos}

\begingroup
\small
\begin{longtable}{p{3.2cm}p{6.2cm}p{4.8cm}}
\toprule
Parte física & Arquivo/função principal & Como executar\\
\midrule
\endhead
Métrica de Kerr &
\texttt{src/kerr\_metric.cpp}, classe \texttt{KerrMetric} &
\texttt{make tests}\\
Geodésicas nulas &
\texttt{src/kerr\_geodesic.cpp}, \texttt{rhs}, \texttt{step\_adaptive} &
\texttt{make kerr-geodesics}\\
Câmera relativística &
\texttt{src/kerr\_camera.cpp}, \texttt{initial\_state}, \texttt{trace\_pixel} &
\texttt{make kerr-rays} ou \texttt{make kerr-geodesics}\\
Toro do collapsar &
\texttt{src/torus\_profile.cpp}, \texttt{rho}, \texttt{Ye}, \texttt{temperature\_MeV} &
\texttt{make plot-torus}\\
Seção de choque DIS &
\texttt{src/sigma\_table.cpp}, \texttt{sigma\_cm2} &
Usado por \texttt{make image-from-cache}\\
Profundidade óptica &
\texttt{src/radiative\_transfer.cpp}, \texttt{integrate\_kerr\_ray} &
\texttt{make image-from-cache ENU=...}\\
Emissividade UHE &
\texttt{src/radiative\_transfer.cpp}, \texttt{emissivity\_collapsar\_ring} &
\texttt{make image-from-cache SOURCE\_R\_RG=...}\\
Canal térmico MeV &
\texttt{src/radiative\_transfer.cpp}, \texttt{emissivity\_mev\_thermal},
\texttt{opacity\_mev\_*} &
\texttt{make image-from-cache MEV\_ENU=10}\\
Imagem final &
\texttt{apps/compute\_kerr\_image\_from\_cache.cpp} &
\texttt{make kerr-criar-cache}\\
\bottomrule
\end{longtable}
\endgroup

\section{Observações e limitações do modelo atual}

\begin{itemize}
  \item O fluxo Schwarzschild foi desativado no \texttt{Makefile}; o projeto
  atual usa Kerr como caminho principal.
  \item O redshift armazenado na câmera é uma aproximação escalar inspirada
  no fator de Schwarzschild, embora os raios sejam propagados na métrica de
  Kerr.
  \item O toro é paramétrico, não uma solução hidrodinâmica autoconsistente.
  \item A seção de choque DIS é tabelada; fora do intervalo da tabela, a
  contribuição é ignorada na transferência radiativa principal.
  \item O canal MeV é uma aproximação física local, sem tabela microfísica e
  sem resolver transporte multi-flavor de Boltzmann; as constantes são
  parametrizações de ordem de grandeza.
  \item A normalização da emissividade é arbitrária; as imagens são mais
  diretamente interpretáveis em termos relativos, contraste e morfologia.
\end{itemize}

\begin{thebibliography}{9}

\bibitem{MacFadyenWoosley1999}
A.~I. MacFadyen and S.~E. Woosley,
``Collapsars: Gamma-Ray Bursts and Explosions in `Failed Supernovae',''
\emph{The Astrophysical Journal} \textbf{524}, 262--289 (1999).
doi: \href{https://doi.org/10.1086/307790}{10.1086/307790}.

\bibitem{MacFadyenWoosley1999Web}
A.~I. MacFadyen,
``Accretion disk density structure at 7.6 s,''
collapsar simulation figure accompanying MacFadyen and Woosley (1999).
\url{https://www.ucolick.org/~andrew/collapsar/fig8.html}.

\bibitem{LiuGuZhang2017}
T.~Liu, W.-M. Gu and B.~Zhang,
``Neutrino-dominated accretion flows as the central engine of gamma-ray bursts,''
\emph{New Astronomy Reviews} \textbf{79}, 1--25 (2017).
doi: \href{https://doi.org/10.1016/j.newar.2017.07.001}
{10.1016/j.newar.2017.07.001}.

\bibitem{ZenatiSiegelMetzger2020}
Y.~Zenati, D.~M. Siegel, B.~D. Metzger and H.~B. Perets,
``Nuclear burning in collapsar accretion discs,''
\emph{Monthly Notices of the Royal Astronomical Society}
\textbf{499}, 4097--4113 (2020).
doi: \href{https://doi.org/10.1093/mnras/staa3002}
{10.1093/mnras/staa3002}.

\bibitem{SheLiuXue2022}
J.-Z. She, T.~Liu and L.~Xue,
``Relativistic global solutions of neutrino-dominated accretion flows with
magnetic coupling,''
\emph{Monthly Notices of the Royal Astronomical Society}
\textbf{513}, 3960--3970 (2022).
doi: \href{https://doi.org/10.1093/mnras/stac1154}
{10.1093/mnras/stac1154}.

\end{thebibliography}

\end{document}
